% GENERATED FILE. Edit canonical lessons, then run npm run generate:books.
% canonical-source-hash: fnv1a64:dcae05eeff150ebb
% canonical-lessons: JA-C124-kayui, JA-C124-kusai, JA-C124-mushiatsui, JA-C124-okashii, JA-C124-omoshiroi

\chapter{Itchy, Smelly, Humid, Funny, Interesting}
\label{ch:ja-itchy-smelly-humid-funny-interesting}
\hlchaptermodality{124}

\begin{chapteropening}
\textbf{By the end of this chapter:} \emph{I can say itchy, smelly, humid, funny and interesting.}

Complete the last lesson of chapter 124: I can say itchy, smelly, humid, funny and interesting.
\end{chapteropening}

\section[\texorpdfstring{\ja{かゆい} (kayui)}{かゆい (kayui)}]{\ja{かゆい} (kayui) — itchy}

\label{lesson:JA-C124-kayui}



\begin{quote}
Before the new one: say the Japanese for suspicious, then the Japanese for detailed.
\end{quote}

\subsection*{You'll want to know: \ja{かゆい}}
\textbf{\ja{かゆい}} — \emph{kayui} — \textquotedblleft{}itchy\textquotedblright{}.

\begin{grammarlens}[title={What you've built}]
One more word for this chapter.
\end{grammarlens}

\subsection*{Your turn}
\begin{itemize}
\raggedright
  \item \emph{Say it:} \emph{kayui}
  \item \emph{Say it:} \emph{kayui}, once more
  \item \emph{Say it:} say \emph{kayui}
  \item \emph{Recall:} say the Japanese for suspicious, then the Japanese for detailed, then say \emph{kayui} again
\end{itemize}

\begin{tcolorbox}[breakable,colback=teal!4,colframe=teal!35!black,title={Before you move on}]
What does \ja{かゆい} mean? (\textquotedblleft{}Itchy\textquotedblright{}.) Say it once more.
\end{tcolorbox}
\section[\texorpdfstring{\ja{くさい} (kusai)}{くさい (kusai)}]{\ja{くさい} (kusai) — smelly}

\label{lesson:JA-C124-kusai}



\begin{quote}
Before the new one: say the Japanese for detailed, then the Japanese for itchy.
\end{quote}

\subsection*{You'll want to know: \ja{くさい}}
\textbf{\ja{くさい}} — \emph{kusai} — \textquotedblleft{}smelly\textquotedblright{}.

\begin{grammarlens}[title={What you've built}]
One more word for this chapter.
\end{grammarlens}

\subsection*{Your turn}
\begin{itemize}
\raggedright
  \item \emph{Say it:} \emph{kusai}
  \item \emph{Say it:} \emph{kusai}, once more
  \item \emph{Say it:} say \emph{kusai}
  \item \emph{Recall:} say the Japanese for detailed, then the Japanese for itchy, then say \emph{kusai} again
\end{itemize}

\begin{tcolorbox}[breakable,colback=teal!4,colframe=teal!35!black,title={Before you move on}]
What does \ja{くさい} mean? (\textquotedblleft{}Smelly\textquotedblright{}.) Say it once more.
\end{tcolorbox}
\section[\texorpdfstring{\ja{むしあつい} (mushiatsui)}{むしあつい (mushiatsui)}]{\ja{むしあつい} (mushiatsui) — humid}

\label{lesson:JA-C124-mushiatsui}



\begin{quote}
Before the new one: say the Japanese for itchy, then the Japanese for smelly.
\end{quote}

\subsection*{You'll want to know: \ja{むしあつい}}
\textbf{\ja{むしあつい}} — \emph{mushiatsui} — \textquotedblleft{}humid\textquotedblright{}.

\begin{grammarlens}[title={What you've built}]
One more word for this chapter.
\end{grammarlens}

\subsection*{Your turn}
\begin{itemize}
\raggedright
  \item \emph{Say it:} \emph{mushiatsui}
  \item \emph{Say it:} \emph{mushiatsui}, once more
  \item \emph{Say it:} say \emph{mushiatsui}
  \item \emph{Recall:} say the Japanese for itchy, then the Japanese for smelly, then say \emph{mushiatsui} again
\end{itemize}

\begin{tcolorbox}[breakable,colback=teal!4,colframe=teal!35!black,title={Before you move on}]
What does \ja{むしあつい} mean? (\textquotedblleft{}Humid\textquotedblright{}.) Say it once more.
\end{tcolorbox}
\section[\texorpdfstring{\ja{おかしい} (okashii)}{おかしい (okashii)}]{\ja{おかしい} (okashii) — funny; strange}

\label{lesson:JA-C124-okashii}



\begin{quote}
Before the new one: say the Japanese for smelly, then the Japanese for humid.
\end{quote}

\subsection*{You'll want to know: \ja{おかしい}}
\textbf{\ja{おかしい}} — \emph{okashii} — \textquotedblleft{}funny; strange\textquotedblright{}.

\begin{grammarlens}[title={What you've built}]
One more word for this chapter.
\end{grammarlens}

\subsection*{Your turn}
\begin{itemize}
\raggedright
  \item \emph{Say it:} \emph{okashii}
  \item \emph{Say it:} \emph{okashii}, once more
  \item \emph{Say it:} say \emph{okashii}
  \item \emph{Recall:} say the Japanese for smelly, then the Japanese for humid, then say \emph{okashii} again
\end{itemize}

\begin{tcolorbox}[breakable,colback=teal!4,colframe=teal!35!black,title={Before you move on}]
What does \ja{おかしい} mean? (\textquotedblleft{}Funny; strange\textquotedblright{}.) Say it once more.
\end{tcolorbox}
\section[\texorpdfstring{\ja{おもしろい} (omoshiroi)}{おもしろい (omoshiroi)}]{\ja{おもしろい} (omoshiroi) — interesting}

\label{lesson:JA-C124-omoshiroi}



\begin{quote}
Before the new one: say the Japanese for humid, then the Japanese for funny; strange.
\end{quote}

\subsection*{You'll want to know: \ja{おもしろい}}
\textbf{\ja{おもしろい}} — \emph{omoshiroi} — \textquotedblleft{}interesting\textquotedblright{}.

\begin{grammarlens}[title={What you've built}]
One more word for this chapter.
\end{grammarlens}

\subsection*{Your turn}
\begin{itemize}
\raggedright
  \item \emph{Say it:} \emph{omoshiroi}
  \item \emph{Say it:} \emph{omoshiroi}, once more
  \item \emph{Say it:} say \emph{omoshiroi}
  \item \emph{Recall:} say the Japanese for humid, then the Japanese for funny; strange, then say \emph{omoshiroi} again
\end{itemize}

\begin{tcolorbox}[breakable,colback=teal!4,colframe=teal!35!black,title={Before you move on}]
What does \ja{おもしろい} mean? (\textquotedblleft{}Interesting\textquotedblright{}.) Say it once more.
\end{tcolorbox}
